\documentclass[10pt]{article}

\usepackage[utf8]{inputenc}

\usepackage[T1]{fontenc}

\usepackage{amsmath}

\usepackage{amsfonts}

\usepackage{amssymb}

\usepackage[version=4]{mhchem}

\usepackage{stmaryrd}

\usepackage{mathrsfs}

\usepackage{bbold}



\begin{document}

В более типичных ситуациях топология $\mathscr{F}$ может быть неотделимой, даже если отделимо $X$ (таков, напр., пучок ростков непрерывных (или дифференцируемых) Функций, порожденный предпучком $F$, где $F(U)-$ непрерывные (диффференцируемые) функции на $U$; однако пучок ростков аналитич. функций на многообразии отделим).



Всякий гомоморфизм предпучков $F \rightarrow F^{\prime}$ приводит к отображению соответствующих пучков $\mathscr{F} \rightarrow \mathscr{F}^{\prime}$, к-рое является локальным гомеоморфизмом и гомоморфно отображает слои в слои; такое отображение пучков наз. г о м о м о р ф и з м о м п у ч к о в. Стандартным образом определяются моно- и эпиморфизмы. При любом гомоморфизме $f: \mathscr{F}^{\prime} \rightarrow \mathscr{J}$ образ $f\left(\mathscr{F}^{\prime}\right)$ есть открытая часть $\mathscr{F}$, замкнутая по отношению к послойным алгебраич. операциям. Всякая часть $\mathscr{F}$, удовлетворяюпая этим требованиям, наз. п о д п у ч к о м в $\mathscr{F}$. Ф а к т о р п у ч о к пучка $\mathscr{F}$ по подпучку $\mathscr{F}^{\prime}$ определяется как пучок $\mathscr{F}^{\prime \prime}$, порожденный предпучком $U \rightarrow \boldsymbol{\Gamma}(U, \mathscr{F}) / \boldsymbol{\Gamma}\left(U, \mathscr{F}^{\prime}\right) ;$ при этом имеется эпиморфизм $\mathscr{F} \rightarrow \mathscr{F}^{\prime \prime}$, причем $\mathscr{F}^{\prime \prime} x^{\prime}=\mathscr{F}_{x} / \mathscr{F}^{\prime} x$. Для всякого открытого $U \subset X$ через $\mathscr{F}_{U}$ обозначается подпучок в $\mathscr{F}$, являющийся объединением $p^{-1}(U)$ с нулевым сечением $\mathscr{F}$ над $X$, а через $\mathscr{F}_{X \backslash U}-$ соответствуюгцй факторпучок (ограничение к-рого на $X \backslash U$ совпадает с ограничением $\mathscr{F})$.



Возможность употреблять по отношению к пучкам над $X$ такие привычные термины, как гомоморфизм, ядро, образ, подпучок, факторпучок и т. д., вкладывая в эти понятия такой же смысл, как в алгебре, позволяет рассматривать их с категорной точки зрения и применять В П. Т. конструкции гомологической алгебры. Возникающие над $X$ категории пучков родственны таким классическим, как категория абелевых групп или категория модулей; в частности, для пучков определяются прямые суммы, бесконечные прямые произведения, индуктивные пределы и др. понятия.



Аппарат П. т. проник в разнообразные области математики благодаря тому, что определены естественные когомологии $H^{*}(X, \mathscr{F})$ пространства $X$ с коәффициентами в пучке $\mathscr{F}$, причем без каких-либо ограничений на $X$ (что суцественно, напр., в алгебраич. геометрии, где возникающие пространства, как правило, неотделимы), и что другие когомологии (в тех или иных конкретных условиях) сводятся к пучковым по крайней мере в тех ситуациях, где их применение оправдано. Для определения $H^{*}(X, \mathscr{F})$ сначала строится к а н оническая ре воль в ента



\[

C^{*}(\mathscr{F}): 0 \longrightarrow \mathscr{\mathscr { T }} \longrightarrow C^{0}(\mathscr{F}) \longrightarrow C^{\mathbf{1}}(\mathscr{F}) \longrightarrow \ldots,

\]



где $C^{0}(\mathscr{F})$ - пучок, определяемый предпучком $F$, для к-рого $F(U)$ - группа всех (включая разрывные) сечений $\mathscr{F}$ над $U$, при этом $\Gamma\left(U, C^{0}(\mathcal{F})\right)=$ $=F(U), \quad C^{1}(\mathscr{F})=C^{0}\left(C^{0}(\mathscr{F}) / \mathscr{F}\right), \ldots, \quad C^{p+1}(\mathscr{F})=$ $=C^{0}\left(C^{p}(\mathscr{F}) / \operatorname{Im} C^{p-1}(\mathscr{F})\right), \cdots$ По огределению, $H^{p}(X, \mathscr{F})=H^{P}\left(\Gamma\left(X, C^{*}(\mathscr{F})\right)\left(H_{\Phi}^{p}(X, \quad \mathscr{F})\right.\right.$ получаются заменой символа $\Gamma$ на $\left.\left.\Gamma_{\Phi}\right)\right)$. При этом сам пучок $\mathscr{F}$ удаляется из $C^{*}(\mathscr{F})$, так что $H_{\Phi}^{0}(X, \mathscr{F})=\Gamma_{\Phi}(X, \mathscr{F})$ (для классич. когомологий $H^{0}(X, G)-$ группа локально постоянных функций на $X$ со значениями в $G)$. Резольвента $C^{*}(\mathscr{F})$-точный ковариантный функтор от $\mathscr{F}$ : точной тройке «коэффициентов» $0 \rightarrow \mathscr{J}^{\prime} \rightarrow \mathscr{F} \rightarrow \mathscr{J}^{\prime \prime} \rightarrow 0$ отвечает точная тройка резольвент. Функтор $\Gamma_{\Phi}$ оказывается точным на членах $C^{p}, p \geqslant 0$, резольвент, поэтому указанным коэфффициентам отвечает точная последовательность когомологий



\[

\begin{gathered}

\cdots \longrightarrow H_{\Phi}^{p-1}\left(X, \mathscr{F}^{\prime \prime}\right) \rightarrow H_{\Phi}^{p}\left(X, \mathscr{F}^{\prime g}\right) \rightarrow \\

\longrightarrow H_{\Phi}^{p}(X, \mathscr{F}) \rightarrow H_{\Phi}^{p}\left(X, \mathscr{F}^{\prime \prime}\right) \longrightarrow \ldots,

\end{gathered}

\]



начинающаяся с $0 \rightarrow \Gamma_{\Phi}\left(X, \mathscr{F}^{\prime}\right) \rightarrow \Gamma_{\Phi}(X, \mathscr{F}) \rightarrow \ldots$. Когомологич. последовательность пары $(X, A)$ отвечает тройке $0 \rightarrow \mathscr{F}_{X \backslash A} \rightarrow \mathscr{F} \rightarrow \mathscr{F}_{A} \rightarrow 0(A-$ замкнутое множество).



Когомологии $H_{\Phi}^{*}(X, \mathscr{F})$ обладают следующим свойством "универсальности", раскрывающим их значение: для любой другой резольвенты $\mathscr{L}^{*}$ (т. е. начинающейся с $\mathscr{F}$ точной последовательности пучков $\mathscr{L}^{q}$ ) имеется естественный гомоморфизм «сравнения» $H^{p}\left(\Gamma_{\Phi}(X\right.$, $\left.\left.\mathscr{L}^{*}\right)\right) \rightarrow H_{\Phi}^{p}(X, \mathscr{F})$, для описания к-рого в терминах $H_{\Phi}^{p}\left(X, \mathscr{L}^{q}\right)$ применяются спектральные последовательности. Важен случай, когда пучки резольвенты Фа ц и к л и ч ны, т. е. когда $H_{\Phi}^{p}\left(X, \mathscr{L}^{q}\right)=0$ при $p \geqslant 1$ : в этом случае указанный гомоморфизм есть изоморфиим. Основными примерами ацикличных пучков являются вялые пучки (для всех $U \in X$ отображения $\Gamma(X, \mathscr{L}) \rightarrow$ $\rightarrow \Gamma(U, \mathscr{L})$ эпиморфны) и мягкие пучки (любое сечение над замкнутым множеством продолжается до сөчения над всем $X$ ). Канонич. резольвента состоит из вялы пучков. Если $X$ - паракомпактное пространство, то всякий вялый пучок является также и мягким.



Свойство универсальности повволяет сравнивать с пучковыми (а следовательно, и между собой) когомологии, возникающие в более конкретных ситуациях, определять для них те естественные границы, в к-рых их применение әффективно, а также применять методы П. т. для решения конкретных задач. Напр., когомологии Александрова - Чеха можно определить с помощью коцепей, получающихся из коцепей специально подобранной системы открытых покрытий переходом к прямому пределу. Эти коцепи оказываются сечениями пучков ростков коцепей (определяемых аналогично пучкам ростков функций), составляющих резольвенту группы (или даже пучка) коәфффициентов, к-рая оказывается мягкой, если пространство паракомпактно. Таким образом, для паракомпактных пространств когомологии Александрова - Чеха совпадают с пучковыми. Аналогичный вывод имеет место для пространств Зариского (в частности, для алгебраич. многообразий). Сечениями пучков резольвенты оказываются и коцепи Александера - Спеньера, причем резольвента состоит из мягких пучков, если $X$ паракомпактно, в частности, в этом случае когомологии Александера - Спеньера и Александрова - Чеха естественно изоморфны. В случае сингулярных когомологий отождествление коцепей, совпадающих друг с другом на сингулярных симплексах мелкости (произвольных) открытых покрытий, приводит к т. н. локализованным коцепям (даюгим те же когомологии), к-рые являются сечениями пучков, определяемых предпучками обычных сингулярных коцепей. Пучки оказываются мягкими, если $X$ паракомпактно (а если $X$ наследственно паракомпактно, то даже вялыми), но образуют резольвенту при дополнительном требовании, чтобы $X$ было слабо локально стягиваемым (в каждой окрестности $U$ каждой точки $x \in X$ найдется менышая окрестность, стягиваемая в точку внутри $U)$. Классич. примером является т е о р е м а д е Р а м а: когомологии комплекса диффференциальных форм дифференцируемого многообразия совпадают с обычными когомологиями с кођффициентами в поле $\mathbb{R}$ действительных чисел (пучки ростков диффференциальных форм являются мягкими и образуют резольвенту $\mathbb{R}:$ вблизи каждой точки каждая замкнутая диффференциальная форма является точной).



Имеются также резольвенты, отвечающие любым открытым или локально конечным замкнутым покрытиям и позволяющие сравнивать когомологии $X$ с когомологиями покрытий (спектральные последовательности покрытий). В частности, изоморфизм обеспечивается





\end{document}